\documentclass[11pt]{article}

\usepackage{amsmath}   % math formulas
\usepackage{amssymb}   % math symbols
\usepackage{graphicx}  % images (not used here, just loaded)
\usepackage{hyperref}  % hyperlinks and cross references

\title{Local \LaTeX{} Environment Test}
\author{Yingyou Ma}
\date{\today}

\begin{document}

\maketitle

\section{Introduction}
Hello! This is a test document for verifying the local VS Code + MiKTeX + latexmk
setup. If you can see this text and the formulas and table below are typeset
correctly, the whole pipeline (compile $\to$ PDF preview) is working.

\section{Math Formulas}
Inline formulas work fine, e.g. Einstein's $E = mc^2$ and Euler's identity
$e^{i\pi} + 1 = 0$.

Display equation with automatic numbering:
\begin{equation}
    \int_0^\infty e^{-x^2}\, dx = \frac{\sqrt{\pi}}{2}.
    \label{eq:gauss}
\end{equation}

A system of aligned equations:
\begin{align}
    \nabla \cdot \mathbf{E} &= \frac{\rho}{\varepsilon_0}, \\
    \nabla \times \mathbf{B} &= \mu_0 \mathbf{J} + \mu_0 \varepsilon_0 \frac{\partial \mathbf{E}}{\partial t}.
\end{align}

We can reference the earlier equation~\eqref{eq:gauss}; cross references work too.

\section{Lists}
Unordered list:
\begin{itemize}
    \item Compiler: MiKTeX
    \item Scheduler: latexmk
    \item Dependency: Strawberry Perl
\end{itemize}

Ordered list:
\begin{enumerate}
    \item Save the file (\texttt{Ctrl+S}) to trigger auto build
    \item The PDF opens in the preview
    \item Double-click the PDF to jump back to the source
\end{enumerate}

\section{Table}
\begin{table}[h]
    \centering
    \begin{tabular}{lcc}
        \hline
        Component & Installed & Role \\
        \hline
        MiKTeX          & Yes & Compile \\
        LaTeX Workshop  & Yes & VS Code extension \\
        Strawberry Perl & Yes & Run latexmk \\
        \hline
    \end{tabular}
    \caption{Overview of the local \LaTeX{} environment components}
    \label{tab:components}
\end{table}

As shown in Table~\ref{tab:components}, all components are in place.

\section{Conclusion}
Everything works! You can now start writing \LaTeX{} locally.

\end{document}
